\documentclass[a4paper,11pt]{article}

\usepackage[utf8]{inputenc}
\usepackage[T1]{fontenc}
\usepackage[french]{babel}
\usepackage{amsmath, amssymb, amsfonts, amsthm}
\usepackage{graphicx}
\usepackage{hyperref}
\usepackage{geometry}
\usepackage{booktabs}
\usepackage{float}
\usepackage{graphicx}
\usepackage{subcaption}
\usepackage{fancyhdr}
\usepackage{xcolor}
\usepackage{listings}
\usepackage{array}
\usepackage{multirow}
\usepackage{enumitem}

\lstset{
  language=Python,
  basicstyle=\small\ttfamily,
  keywordstyle=\color{blue},
  commentstyle=\color{gray},
  stringstyle=\color{orange},
  frame=single,
  breaklines=true,
  numbers=left,
  numberstyle=\tiny\color{gray},
  xleftmargin=1.5em,
  showspaces=false,
  showstringspaces=false
}

\pagestyle{fancy}
\renewcommand{\headrulewidth}{0pt}
\fancyhead[C]{%
    \makebox[\textwidth][c]{Data Battle 2026 --- Rapport de Modélisation Probabiliste}\\
    \makebox[\textwidth][c]{\rule{0.9\textwidth}{0.1pt}}
}
\fancyhead[L]{}
\fancyhead[R]{}
\fancyfoot[L]{}
\fancyfoot[R]{}
\fancyfoot[C]{\makebox[\textwidth][c]{\thepage}}

\geometry{hmargin=2cm,vmargin=2.5cm}

\newtheorem{definition}{Définition}[section]
\newtheorem{proposition}{Proposition}[section]
\newtheorem{remarque}{Remarque}[section]
\newtheorem{theoreme}{Théorème}[section]

\title{\textbf{Estimation Probabiliste de la Fin d'Alerte Orage\\
\large Modélisation par Processus de Hawkes et Gradient Boosting}}
\author{}
\date{\textit{Data Battle 2026 --- Meteorage}}

\begin{document}

\maketitle

\begin{abstract}
Ce rapport présente une approche en deux volets pour estimer la probabilité de fin d'une alerte foudre autour d'aéroports européens, à partir de données d'impacts fournis par Meteorage. Le premier volet, probabiliste, repose sur la modélisation des impacts comme un processus ponctuel auto-excitant de Hawkes, dont les paramètres sont estimés par maximum de vraisemblance, et dont la validité est évaluée par le théorème de rescaling temporel et par simulation Monte Carlo. À partir du modèle calibré, on calcule en temps réel un délai optimal de levée d'alerte $u^*$ sous un risque toléré $\varepsilon$. Le second volet, supervisé, traduit le même problème en une tâche de régression sur des features temporelles construites à partir des historiques d'impacts, et utilise LightGBM pour apprendre les patterns de fin d'alerte directement depuis les données. Les deux approches sont complémentaires : la première fournit une interprétabilité théorique et une quantification rigoureuse de l'incertitude, la seconde offre une flexibilité et une puissance prédictive supérieures lorsque le volume de données le permet.
\end{abstract}

\tableofcontents
\newpage

%=============================================================
\section{Vue d'ensemble des deux stratégies}
%=============================================================

Le tableau~\ref{tab:comparaison} résume les deux approches développées dans ce rapport et leur articulation.

\begin{table}[H]
\centering
\renewcommand{\arraystretch}{1.4}
\begin{tabular}{>{\bfseries}p{3.2cm} p{5.8cm} p{5.8cm}}
\toprule
 & \textbf{Approche 1 : Processus de Hawkes} & \textbf{Approche 2 : LightGBM} \\
\midrule
Paradigme & Modèle génératif probabiliste & Modèle discriminatif supervisé \\
Objectif direct & Estimer $\mathbb{P}(T > \tau_n + u \mid \mathcal{H}_n)$ & Prédire $u^*$ à partir de features \\
Entrées & Séquences temporelles $\{\tau_i\}$ & Vecteur de features $\mathbf{f}_n$ construit depuis $\mathcal{H}_n$ \\
Paramètres & $(\hat\mu, \hat\alpha, \hat\beta)$ par MLE & Arbres de décision boostés, appris par CV \\
Validation & Théorème de rescaling + simulation Monte Carlo & Métriques de régression (MAE, RMSE) sur ensemble de test \\
Interprétabilité & Totale : chaque paramètre a un sens physique & Partielle : via SHAP values \\
Hypothèses & Stationnarité par aéroport, noyau exponentiel & Aucune hypothèse distributionnelle \\
Avantages & Quantification de l'incertitude, extrapolation & Capture les non-linéarités, les marques \\
Limites & Hypothèse paramétrique forte, $\sim$150 alertes & Boîte noire, pas de distribution de $u^*$ \\
\bottomrule
\end{tabular}
\caption{Comparaison des deux stratégies de modélisation.}
\label{tab:comparaison}
\end{table}

\begin{remarque}
Nous recommandons de développer les deux approches en parallèle. L'approche probabiliste est indispensable pour la rigueur mathématique, la quantification de l'incertitude et la compréhension du phénomène. LightGBM sera plus compétitif sur le score final car il peut intégrer naturellement les marques (amplitude, position, type d'éclair) sans hypothèse de forme sur leur rôle dans la dynamique.
\end{remarque}

%=============================================================
\section{Formulation mathématique du problème}
%=============================================================

\subsection{Espace de probabilité et données}

On se place sur un espace de probabilité filtré $(\Omega, \mathcal{F}, (\mathcal{F}_t)_{t \geq 0}, \mathbb{P})$. Soit $A \in \mathbb{R}^2$ le centre géographique d'un aéroport et $r = 20$ km. On note :
\[
\mathcal{D} = B(A, r) = \{ x \in \mathbb{R}^2 \mid \|x - A\| \leq r \}
\]
le disque de surveillance dans lequel les impacts déclenchent et prolongent les alertes.

\subsection{Processus ponctuel marqué}

Les impacts de foudre survenant dans $\mathcal{D}$ sont modélisés par un \textbf{processus ponctuel marqué} :
\[
\mathcal{N} = \left\{ \left( \tau_i,\, \mathbf{X}_i \right) \right\}_{i \geq 1}
\]
où, pour chaque impact $i$ :
\begin{itemize}
    \item $\tau_i \in \mathbb{R}_+$ est l'\textbf{instant d'impact} (colonne \texttt{date}),
    \item $\mathbf{X}_i = (\xi_i,\, \phi_i,\, \alpha_i,\, d_i,\, C_i) \in \mathcal{X}$ est le \textbf{vecteur de marques} composé de :
    \begin{itemize}
        \item $(\xi_i, \phi_i) \in \mathcal{D}$ : position géographique (\texttt{lon}, \texttt{lat}),
        \item $\alpha_i \in \mathbb{R}$ : amplitude du signal en kA (\texttt{amplitude}),
        \item $d_i \in [0, r]$ : distance à l'aéroport (\texttt{dist}),
        \item $C_i \in \{0, 1\}$ : type d'éclair, nuage-sol ou intra-nuageux (\texttt{icloud}).
    \end{itemize}
\end{itemize}

\subsection{Définition formelle d'une alerte}

\begin{definition}[Alerte]
Une \textbf{alerte} est déclenchée à l'instant $\tau_0$ dès qu'un impact survient dans $\mathcal{D}$. Elle se termine à la première période de $\Delta = 30$ minutes consécutives sans impact dans $\mathcal{D}$. La \textbf{durée totale de l'alerte} est la variable aléatoire :
\[
T = \inf \left\{ t > \tau_0 \;\middle|\; \mathcal{N}\!\left((t,\, t + \Delta) \cap \mathcal{D} \times \mathcal{X}\right) = 0 \right\} + \Delta - \tau_0.
\]
\end{definition}

\subsection{Historique et filtration}

On se place à l'instant courant $\tau_n$, instant du $n$-ième impact dans $\mathcal{D}$, avec :
\[
\tau_0 < \tau_1 < \cdots < \tau_n, \qquad \tau_i - \tau_{i-1} < \Delta \quad \forall\, i \in \{1, \ldots, n\}.
\]
On définit la \textbf{filtration naturelle} du processus et l'\textbf{historique observé} jusqu'à $\tau_n$ :
\[
\mathcal{H}_n = \sigma\!\left( \tau_i,\, \xi_i, \phi_i, \alpha_i, d_i, C_i \;\middle|\; 0 \leq i \leq n \right).
\]
C'est ce vecteur aléatoire observé qui constitue la statistique exhaustive permettant toute inférence sur la dynamique future du processus.

\subsection{Objectif probabiliste}

\begin{definition}[Délai optimal de levée d'alerte]
Soit $\varepsilon \in (0,1)$ le \textbf{niveau de risque toléré} par l'aéroport. À chaque nouvel impact en $\tau_n$, on cherche le \textbf{délai minimal de levée d'alerte} $u^* \in [0, \Delta]$ tel que :
\[
\boxed{u^* = \inf \left\{ u \in [0, \Delta] \;\middle|\; \mathbb{P}\!\left( T > \tau_n + u \;\middle|\; \mathcal{H}_n \right) < \varepsilon \right\}.}
\]
\end{definition}

\textbf{Interprétation.} Après $u^*$ minutes sans nouvel éclair depuis $\tau_n$, on peut admettre sous risque $\varepsilon$ que l'alerte est terminée, permettant à l'aéroport de reprendre son activité avant l'expiration du délai réglementaire fixe de $\Delta = 30$ min.

%=============================================================
\section{Stratégie 1 : Modélisation par processus de Hawkes}
%=============================================================

\subsection{Choix du modèle}

Pour modéliser la dynamique temporelle des impacts dans $\mathcal{D}$, nous choisissons un \textbf{processus de Hawkes univarié à noyau exponentiel}. Ce choix est motivé par la propriété d'auto-excitation : un impact augmente la probabilité d'impacts proches dans le temps, ce qui est physiquement cohérent avec la dynamique des orages où un éclair favorise la création de nouveaux éclairs.

\begin{definition}[Processus de Hawkes à noyau exponentiel]
L'intensité conditionnelle du processus est définie par :
\[
\lambda^*(t) = \mu + \alpha \sum_{\tau_i < t} e^{-\beta(t - \tau_i)},
\]
où :
\begin{itemize}
    \item $\mu > 0$ est le \textbf{taux de base} (activité orageuse de fond),
    \item $\alpha > 0$ est le paramètre d'\textbf{excitabilité} (amplitude du saut à chaque impact),
    \item $\beta > 0$ est le \textbf{taux de décroissance} (vitesse d'atténuation de l'excitation).
\end{itemize}
\end{definition}

\begin{remarque}
La condition de stationnarité du processus exige $\alpha < \beta$, soit $\alpha/\beta < 1$. Ce ratio mesure la \textbf{branchabilité} du processus : un ratio proche de 1 indique un processus proche de la criticité, où chaque impact génère presque autant d'impacts secondaires qu'il n'en reçoit. Les résultats obtenus (ratio $\approx 0.98$ pour tous les aéroports) indiquent des orages très auto-excités, cohérents avec la physique orageuse.
\end{remarque}

\subsection{Lien avec la probabilité de survie}

Pour établir l'expression de $\mathbb{P}(T > \tau_n + u \mid \mathcal{H}_n)$,
on procède en deux étapes : d'abord sur un processus de Poisson
\emph{déterministe} pour faire apparaître l'exponentielle, puis on
revient au cas du Hawkes où l'intensité est \emph{aléatoire} pour
faire apparaître l'espérance.

\subsubsection*{Étape 1 — Processus de Poisson inhomogène à intensité
déterministe}

Supposons provisoirement que $\lambda^*(s) = \lambda(s)$ est une
fonction \textbf{déterministe} du temps (non aléatoire). Le nombre
d'impacts $N([a,b])$ dans un intervalle $[a,b]$ suit alors une loi de
Poisson de paramètre :
\[
    \Lambda(a,b) = \int_a^b \lambda(s)\,ds.
\]
La probabilité qu'il y en ait \textbf{zéro} est donnée par la formule
de la loi de Poisson pour $k = 0$ :
\[
    \mathbb{P}\!\bigl(N([a,b]) = 0\bigr)
    = \frac{e^{-\Lambda(a,b)} \Lambda(a,b)^0}{0!}
    = e^{-\Lambda(a,b)}
    = \exp\!\left(-\int_a^b \lambda(s)\,ds\right).
\]
C'est de là que vient l'exponentielle. On note en particulier que
$\Lambda(a,b)$ est un \emph{nombre d'impacts attendus}, pas une
probabilité : il peut dépasser 1, mais $e^{-\Lambda}$ reste bien dans
$[0,1]$.

\subsubsection*{Étape 2 — Retour au processus de Hawkes : l'intensité
est aléatoire pour $s > \tau_n$}

Dans le processus de Hawkes, l'intensité conditionnelle est :
\[
    \lambda^*(s)
    = \mu + \alpha \sum_{\tau_i < s} e^{-\beta(s - \tau_i)}.
\]
À l'instant courant $\tau_n$, la somme porte sur tous les impacts
$\tau_i < s$. Or pour $s > \tau_n$, cette somme inclut les impacts
$\tau_{n+1}, \tau_{n+2}, \ldots$ qui \textbf{pourraient survenir dans
$(\tau_n, s)$} et que l'on ne connaît pas encore. Ainsi, pour
$s > \tau_n$ :
\[
    \lambda^*(s) \text{ est une \textbf{variable aléatoire},
    dépendant de la trajectoire future du processus.}
\]
Par conséquent, $\int_{\tau_n}^{\tau_n+u} \lambda^*(s)\,ds$ est
elle-même aléatoire.

\subsubsection*{Étape 3 — Mise en espérance par conditionnement total}

On note $\mathcal{F}$ la tribu engendrée par la trajectoire complète
du processus (passée \emph{et} future). Par la formule des
probabilités totales, en conditionnant d'abord à $\mathcal{F}$ puis
en intégrant sur les trajectoires futures :
\[
    \mathbb{P}(T > \tau_n + u \mid \mathcal{H}_n)
    = \mathbb{E}\!\Bigl[
        \mathbb{P}\!\bigl(T > \tau_n + u \mid \mathcal{F}\bigr)
      \;\Big|\; \mathcal{H}_n
    \Bigr].
\]
Conditionnellement à $\mathcal{F}$, la trajectoire future est
\emph{entièrement fixée}, donc $\lambda^*(s)$ est déterministe
sur chaque trajectoire $\omega$ et l'étape 1 s'applique :
\[
    \mathbb{P}\!\bigl(T > \tau_n + u \mid \mathcal{F}\bigr)(\omega)
    = \exp\!\left(
        -\int_{\tau_n}^{\tau_n+u} \lambda^*_\omega(s)\,ds
      \right).
\]
En substituant dans la formule des probabilités totales :
\[
\boxed{
    \mathbb{P}\!\left(T > \tau_n + u \;\middle|\; \mathcal{H}_n\right)
    = \mathbb{E}\!\left[
        \exp\!\left(
          -\int_{\tau_n}^{\tau_n+u} \lambda^*(s)\,ds
        \right)
      \Bigg|\, \mathcal{H}_n
    \right].
}
\]

\subsubsection*{Pourquoi cette espérance n'a pas d'expression
analytique}

Si $\lambda^*(s)$ était déterministe (Poisson inhomogène pur), le
calcul serait terminé. Dans le Hawkes, $\lambda^*(s)$ dépend de ses
propres réalisations futures : chaque nouvel impact augmente
$\lambda^*$, ce qui rend probable un impact suivant, qui à son tour
augmente $\lambda^*$, etc. Ce \textbf{couplage auto-excitant} entre
l'intensité et les événements qu'elle génère crée une dépendance
récursive qui empêche le calcul en forme fermée.

On estime donc cette espérance par \textbf{simulation Monte Carlo} :
on simule $N$ trajectoires futures $\omega^{(1)}, \ldots, \omega^{(N)}$
depuis $\tau_n$ à l'aide de l'algorithme de thinning, et on estime :
\[
    \mathbb{P}(T > \tau_n + u \mid \mathcal{H}_n)
    \;\approx\;
    \frac{1}{N}
    \sum_{j=1}^{N}
    \mathbf{1}\!\left[
      \text{aucun impact dans } (\tau_n,\, \tau_n + u)
      \text{ dans la simulation } \omega^{(j)}
    \right],
\]
ce qui revient à compter la proportion de simulations où l'alerte
est encore active $u$ minutes après le dernier impact observé.
\subsection{Variable récursive et calcul en $O(n)$}

Un point technique central est l'introduction de la variable récursive :
\[
R_i = \sum_{j=0}^{i-1} e^{-\beta(\tau_i - \tau_j)}, \quad R_0 = 0,
\]
qui se met à jour en $O(1)$ à chaque nouveau pas de temps :
\[
R_i = e^{-\beta(\tau_i - \tau_{i-1})}\,(1 + R_{i-1}).
\]
L'intensité conditionnelle s'écrit alors $\lambda^*(\tau_i) = \mu + \alpha R_i$, et toutes les sommes sur l'historique se calculent sans boucle imbriquée.

%=============================================================
\section{Prétraitement des données --- \texttt{preprocessing.py}}
%=============================================================

\subsection{Description}

Le module \texttt{preprocessing.py} lit le fichier CSV brut et construit la structure de données fondamentale : les \textbf{trajectoires d'alertes}, unités d'observation de l'inférence statistique.

\subsection{Procédure}

\begin{enumerate}
    \item Chargement du CSV et parsing des dates en timestamps UTC.
    \item Filtrage sur \texttt{dist} $\leq 20$ km : seuls ces impacts appartiennent à $\mathcal{D}$.
    \item Pour chaque aéroport, tri chronologique et \textbf{segmentation en alertes} : dès qu'un silence de $\Delta = 30$ min est détecté entre deux impacts consécutifs, l'alerte courante se termine et une nouvelle débute au prochain impact.
    \item Exclusion des alertes à un seul impact (pas d'inter-arrivée exploitable pour la log-vraisemblance).
\end{enumerate}

\subsection{Interface}

\begin{lstlisting}
def load_data(path: str) -> tuple[pd.DataFrame, pd.DataFrame, dict]:
    """
    Entrees :
        path : str  -- chemin vers le CSV brut (507 071 lignes)

    Sorties :
        df       : DataFrame  -- donnees brutes completes
        df_inner : DataFrame  -- impacts filtres dist <= 20 km
        alerts   : dict       -- {airport: [traj_1, ..., traj_K]}
                                 traj_k = array numpy de timestamps
                                 en secondes relatifs a tau_0=0
    """
\end{lstlisting}

\textbf{Structure de la sortie.} Le dictionnaire \texttt{alerts} a la forme :
\[
\texttt{alerts}[\text{airport}] = \left[ \underbrace{[0,\, \tau_1^{(1)},\, \ldots,\, \tau_{n_1}^{(1)}]}_{\text{alerte 1}},\; \ldots,\; \underbrace{[0,\, \tau_1^{(K)},\, \ldots,\, \tau_{n_K}^{(K)}]}_{\text{alerte }K} \right].
\]
Chaque trajectoire $\omega_k$ est un array numpy de timestamps en secondes depuis $\tau_0^{(k)} = 0$. Le tableau~\ref{tab:alertes} donne la répartition des alertes.

\begin{table}[H]
\centering
\begin{tabular}{lcc}
\toprule
Aéroport & Alertes & Impacts dans $\mathcal{D}$ \\
\midrule
Ajaccio  & 34 & $\sim$800 \\
Bastia   & 30 & $\sim$700 \\
Biarritz & 25 & $\sim$600 \\
Nantes   & 38 & $\sim$900 \\
Pise     & 25 & $\sim$600 \\
\midrule
\textbf{Total} & \textbf{152} & $\sim$\textbf{3600} \\
\bottomrule
\end{tabular}
\caption{Répartition des alertes par aéroport après prétraitement.}
\label{tab:alertes}
\end{table}

%=============================================================
\section{Analyse exploratoire --- \texttt{eda.py}}
%=============================================================

\subsection{Description}

Le module \texttt{eda.py} calcule et visualise les distributions empiriques des trois quantités clés pour justifier le choix du modèle de Hawkes et identifier les différences géographiques.

\subsection{Quantités analysées}

\begin{enumerate}
    \item \textbf{Inter-arrivées} $\delta_i = \tau_i - \tau_{i-1}$ (en minutes) : si elles décroissent exponentiellement, cela est cohérent avec un processus de Poisson inhomogène et compatible avec le Hawkes.
    \item \textbf{Durées d'alertes} $T_k = \tau_{n_k}^{(k)} + \Delta$ (en minutes) : donne la distribution de la variable que l'on cherche à prédire.
    \item \textbf{Nombre d'impacts par alerte} $n_k$ : révèle l'intensité typique des orages par géographie.
\end{enumerate}

\subsection{Interface}

\begin{lstlisting}
def eda(alerts: dict) -> None:
    """
    Entrees :
        alerts : dict  -- sortie de load_data()

    Sorties :
        Figures sauvegardees dans figures/eda_*.png
        (histogrammes superposes par aeroport)
    """
\end{lstlisting}

%=============================================================
\section{Estimation par maximum de vraisemblance --- \texttt{hawkes.py}}
%=============================================================

\subsection{Log-vraisemblance d'Ogata}

Pour une trajectoire $\omega_k = \{\tau_0, \tau_1, \ldots, \tau_n\}$ observée sur $[0, T_k]$ avec $T_k = \tau_n + \Delta$, la log-vraisemblance du processus de Hawkes est donnée par le théorème de Ogata (1988) :

\[
\ell_k(\mu, \alpha, \beta) = \underbrace{\sum_{i=0}^{n} \log \lambda^*(\tau_i)}_{\text{contribution de chaque impact}} - \underbrace{\int_0^{T_k} \lambda^*(s)\,ds}_{\text{terme de compensation}}.
\]

Le terme intégral se calcule analytiquement :
\[
\int_0^{T_k} \lambda^*(s)\,ds = \mu T_k + \frac{\alpha}{\beta} \sum_{i=0}^{n} \left(1 - e^{-\beta(T_k - \tau_i)}\right).
\]

Sur l'ensemble des $K$ alertes d'un aéroport, on maximise la log-vraisemblance totale :
\[
\mathcal{L}(\mu, \alpha, \beta) = \sum_{k=1}^{K} \ell_k(\mu, \alpha, \beta).
\]

\subsection{Estimateur du maximum de vraisemblance}

\begin{definition}[EMV du Hawkes]
L'estimateur du maximum de vraisemblance est :
\[
(\hat\mu, \hat\alpha, \hat\beta) = \arg\max_{\mu > 0,\, \alpha > 0,\, \beta > 0} \;\mathcal{L}(\mu, \alpha, \beta).
\]
La maximisation est réalisée numériquement par l'algorithme L-BFGS-B avec une grille de points d'initialisation pour éviter les minima locaux.
\end{definition}

\subsection{Résultats obtenus}

Le tableau~\ref{tab:params} présente les paramètres estimés pour chaque aéroport.

\begin{table}[H]
\centering
\begin{tabular}{lcccc}
\toprule
Aéroport & $\hat\mu$ & $\hat\alpha$ & $\hat\beta$ & $\hat\alpha/\hat\beta$ \\
\midrule
Ajaccio  & 0.00376 & 4.235 & 4.335 & 0.977 \\
Bastia   & 0.00363 & 4.966 & 5.031 & 0.987 \\
Biarritz & 0.00450 & 5.524 & 5.612 & 0.984 \\
Nantes   & 0.00240 & 5.289 & 5.379 & 0.983 \\
Pise     & 0.00591 & 4.681 & 4.762 & 0.983 \\
\bottomrule
\end{tabular}
\caption{Paramètres du processus de Hawkes estimés par MLE pour chaque aéroport.}
\label{tab:params}
\end{table}

\begin{remarque}
Le ratio $\hat\alpha/\hat\beta$ proche de 1 pour tous les aéroports indique des processus proches de la criticité : chaque impact génère en moyenne $\hat\alpha/\hat\beta \approx 0.98$ impacts secondaires. Pise présente le taux de base $\hat\mu$ le plus élevé (0.00591), cohérent avec une activité orageuse de fond plus forte en Toscane.
\end{remarque}

\subsection{Interface}

\begin{lstlisting}
def neg_log_likelihood(params: list, trajectories: list) -> float:
    """
    Entrees :
        params       : [mu, alpha, beta]  -- parametres courants
        trajectories : list of np.array   -- trajectoires d'une meme ville

    Sortie :
        float  -- valeur de -L(mu, alpha, beta) a minimiser
    """

def fit_hawkes(trajectories: list) -> dict:
    """
    Entrees :
        trajectories : list of np.array  -- alertes d'un aeroport

    Sortie :
        dict avec cles 'mu', 'alpha', 'beta', 'result' (objet scipy)
    """

def fit_all_airports(alerts: dict) -> dict:
    """
    Entrees :
        alerts : dict  -- sortie de load_data()

    Sortie :
        dict {airport: params_dict}  -- un modele calibre par aeroport
    """
\end{lstlisting}

%=============================================================
\section{Validation du modèle --- \texttt{validation.py}}
%=============================================================

\subsection{Théorème de rescaling temporel}

La validation repose sur le résultat fondamental suivant.

\begin{theoreme}[Papangelou, 1972 --- Rescaling temporel]
Soit $\{\tau_i\}$ une réalisation d'un processus ponctuel d'intensité conditionnelle $\lambda^*(t)$. On définit la \textbf{compensatrice} :
\[
\Lambda(\tau_i) = \int_0^{\tau_i} \lambda^*(s)\,ds.
\]
Si le modèle est correctement spécifié, les accroissements :
\[
U_i = \Lambda(\tau_i) - \Lambda(\tau_{i-1}), \quad i = 1, 2, \ldots
\]
sont des variables aléatoires indépendantes et identiquement distribuées de loi $\mathrm{Exp}(1)$.
\end{theoreme}

\subsection{Calcul analytique des résidus}

Pour le noyau exponentiel, l'accroissement de la compensatrice entre $\tau_{i-1}$ et $\tau_i$ se calcule analytiquement :
\[
U_i = \mu(\tau_i - \tau_{i-1}) + \frac{\alpha R_{i-1}}{\beta}\left(1 - e^{-\beta(\tau_i - \tau_{i-1})}\right),
\]
où $R_{i-1} = \sum_{j < i-1} e^{-\beta(\tau_{i-1} - \tau_j)}$ est la variable récursive évaluée en $\tau_{i-1}$.

\subsection{Test statistique}

On applique un \textbf{test de Kolmogorov-Smirnov} sur l'ensemble des résidus $\{U_i\}$ contre la loi $\mathrm{Exp}(1)$, complété d'un QQ-plot pour le diagnostic visuel.

\textbf{Résultat obtenu.} Le test est rejeté à 5\% pour Ajaccio (KS = 0.0798, $p < 0.001$). Ce rejet n'invalide pas l'approche : il indique que le modèle homogène de base peut être amélioré, notamment en intégrant les marques (amplitude, distance) dans l'intensité conditionnelle, ou en ajoutant une composante saisonnière dans $\mu$.

\subsection{Validation par simulation}

La seconde approche de validation, décrite par le module \texttt{simulation\_validation()}, consiste à :
\begin{enumerate}
    \item Simuler $N_{\text{sim}} = 500$ alertes avec les paramètres $(\hat\mu, \hat\alpha, \hat\beta)$ estimés,
    \item Comparer les distributions simulées (durées, nombre d'impacts, inter-arrivées) aux distributions observées par superposition des histogrammes.
\end{enumerate}

Si les distributions se superposent, le modèle capture la dynamique essentielle des données.

\subsection{Interface}

\begin{lstlisting}
def compute_residuals(params: dict, trajectories: list) -> np.ndarray:
    """
    Entrees :
        params       : dict          -- {mu, alpha, beta}
        trajectories : list          -- alertes d'un aeroport

    Sortie :
        np.ndarray  -- array des residus 
                       U_i = Lambda(tau_i) - Lambda(tau_{i-1})
                       doit suivre Exp(1) si le modele est correct
    """

def goodness_of_fit(params: dict, trajectories: list, airport: str):
    """
    Entrees :
        params, trajectories, airport  -- idem ci-dessus

    Sorties :
        (stat, pvalue)  -- statistique KS et p-value
        Figures : QQ-plot et histogramme vs Exp(1) dans figures/
    """

def simulation_validation(params: dict, trajectories: list,
                           airport: str, n_sim: int = 500):
    """
    Entrees :
        n_sim : int  -- nombre d'alertes simulees (defaut 500)

    Sorties :
        Figure comparant distributions observees vs simulees
    """
\end{lstlisting}

%=============================================================
\section{Simulation par algorithme de thinning --- \texttt{simulation.py}}
%=============================================================

\subsection{Algorithme de thinning d'Ogata}

La simulation exacte du processus de Hawkes repose sur l'algorithme de thinning d'Ogata (1981). La difficulté est que l'intensité $\lambda^*(t)$ varie dans le temps, empêchant la simulation directe par inversion.

\begin{description}
    \item[Principe.] On construit une borne supérieure locale $\bar\lambda \geq \lambda^*(t)$ sur un intervalle $[t, t + \delta]$. On génère un candidat $t' = t + E$ où $E \sim \mathrm{Exp}(\bar\lambda)$. On accepte $t'$ comme vrai événement avec probabilité $\lambda^*(t')/\bar\lambda$, sinon on le rejette (c'est le « thinning » des candidats fantômes).
\end{description}

Pour le noyau exponentiel, l'intensité est décroissante entre deux événements, donc $\bar\lambda = \lambda^*(t^+)$ (intensité juste après le dernier événement) est une borne valide. La mise à jour récursive :
\[
R_{\text{new}} = e^{-\beta \cdot dt}(1 + R_{\text{old}})
\]
assure un coût $O(1)$ par candidat.

\subsection{Interface}

\begin{lstlisting}
def simulate_alert(mu: float, alpha: float, beta: float,
                   rng=None) -> np.ndarray:
    """
    Simule une alerte complete depuis tau_0 = 0
    jusqu'a la terminaison naturelle (silence de 30 min).

    Entrees :
        mu, alpha, beta : float  -- parametres du processus
        rng : np.random.Generator (optionnel)

    Sortie :
        np.ndarray  -- timestamps en secondes depuis tau_0 = 0
    """

def simulate_hawkes(mu: float, alpha: float, beta: float,
                    T_max: float, history: list = None,
                    rng=None) -> list:
    """
    Simule le processus sur [0, T_max], eventuellement
    conditionne a un historique passe (history).

    Entrees :
        T_max   : float  -- horizon de simulation (secondes)
        history : list   -- evenements passes (timestamps)

    Sortie :
        list  -- nouveaux evenements simules dans [0, T_max]
    """
\end{lstlisting}

%=============================================================
\section{Calcul du délai optimal $u^*$ --- \texttt{prediction.py}}
%=============================================================

\subsection{Estimation de la courbe de survie par Monte Carlo}

L'estimation de $\mathbb{P}(T > \tau_n + u \mid \mathcal{H}_n)$ est réalisée par simulation Monte Carlo.

\begin{description}
    \item[Initialisation.] On initialise $R$ depuis l'historique observé :
    \[
    R = \sum_{i=0}^{n} e^{-\beta(\tau_n - \tau_i)}.
    \]
    \item[Simulation.] On lance $N_{\text{sim}}$ réalisations de la continuation du processus depuis $\tau_n$, en enregistrant pour chaque simulation le délai $u^{(j)}$ jusqu'à la terminaison de l'alerte.
    \item[Estimation.] La courbe de survie est estimée par :
    \[
    \hat{S}(u) = \mathbb{P}(T > \tau_n + u \mid \mathcal{H}_n) \approx \frac{1}{N_{\text{sim}}} \sum_{j=1}^{N_{\text{sim}}} \mathbf{1}\left[u^{(j)} > u\right].
    \]
    \item[Calcul de $u^*$.] On retourne le premier $u$ de la grille $[0, \Delta]$ tel que $\hat{S}(u) < \varepsilon$.
\end{description}

\subsection{Interface}

\begin{lstlisting}
def estimate_survival_and_ustar(
        params: dict,
        history: np.ndarray,
        epsilon: float = 0.10,
        n_sim: int = 3000,
        rng=None) -> dict:
    """
    Entrees :
        params  : dict       -- {mu, alpha, beta} de l'aeroport
        history : np.ndarray -- timestamps en secondes depuis tau_0,
                                dernier element = tau_n (instant courant)
        epsilon : float      -- risque tolere (ex. 0.10 pour 10%)
        n_sim   : int        -- nombre de simulations Monte Carlo

    Sorties :
        dict avec :
            'u_star'   : float      -- delai optimal en secondes
            'u_grid'   : np.ndarray -- grille de u 
                                       (300 points dans [0, Delta])
            'survival' : np.ndarray -- S(u) estimee sur u_grid
        Figure : courbe de survie avec u* et epsilon dans figures/
    """
\end{lstlisting}

\subsection{Exemple de résultat attendu}

Pour une alerte en cours avec $n = 5$ impacts observés et $\varepsilon = 10\%$, le résultat typique est de la forme :
\[
u^* \approx 18 \text{ min} \quad \Longrightarrow \quad \text{Gain de } 30 - 18 = 12 \text{ min sur l'alerte réglementaire.}
\]

% %=============================================================
% \section{Conclusion}
% %=============================================================

% Ce rapport présente un pipeline complet pour l'estimation probabiliste de la fin d'alerte foudre, articulé autour de deux stratégies complémentaires. La stratégie de Hawkes fournit un cadre mathématique rigoureux avec une quantification explicite de l'incertitude via la courbe de survie conditionnelle $\hat{S}(u)$, validée par le théorème de rescaling temporel et par simulation Monte Carlo. La stratégie LightGBM complète cette approche en apprenant les patterns non paramétriques depuis les données, en intégrant naturellement les marques, et en bénéficiant de la sortie du Hawkes comme feature.

% Le code Python associé est organisé en modules indépendants (\texttt{preprocessing}, \texttt{eda}, \texttt{hawkes}, \texttt{validation}, \texttt{simulation}, \texttt{prediction}, \texttt{features}, \texttt{lightgbm\_model}) permettant de tester et améliorer chaque composante séparément. Les pistes d'amélioration les plus prometteuses sont le Hawkes marqué intégrant amplitude et distance, et la calibration de probabilité du LightGBM pour obtenir des intervalles de confiance sur $u^*$.

\end{document}